\subsection{A bounded finite quotient of the actual Hida trace}
\label{sec:cubic-finite-trace-quotient}

Let \(p>2\), let \(\mathcal O\) be the ring of integers of a finite
extension \(E_0/\mathbb Q_p\), and put
\[
 \Lambda=\mathcal O[[W]],\qquad
 \mathfrak m=(\varpi,W),\qquad b=[E_0:\mathbb Q_p].
\]
Let \(S\) be a normal domain finite free of rank \(D\) over \(\Lambda\),
and write \(K=\operatorname{Frac}(S)\) and
\(K_0=\operatorname{Frac}(\Lambda)\).
All occurrences of continuity below use the \(\mathfrak m\)-adic
topology on finite \(\Lambda\)-modules.
The application is the normalized branch of AC-G78 and AC-G84.
Oddness will supply a model over \(K\) itself, without an additional
splitting field.

\begin{lemma}[AC-G85-split: oddness splits the generic degree-two determinant]
\label{lem:G85-split}
Let \(\Gamma\) be a group and let \(\mathscr D:S[\Gamma]\to S\)
be a determinant of dimension two, with trace \(T\).
Suppose its generic semisimple representation over \(\overline K\)
is absolutely irreducible.
If there is \(c\in\Gamma\) with
\[
 c^2=1,\qquad T(c)=0,\qquad \mathscr D(c)=-1,
 \tag{G85.1}
\]
then this representation has a model
\(\rho:\Gamma\to\operatorname{GL}_2(K)\) with trace \(T\).

The generic irreducibility assumption follows, in particular, if
\(\mathscr D\) has an absolutely irreducible specialization to a
characteristic-zero field.
\end{lemma}

\begin{proof}
The existence of the generic semisimple representation is the
field-valued determinant theorem of Chenevier, Theorem 2.12 in the
author's \emph{determinants.pdf}; this published input is retained.
Since \(c^2=1\) and \(2\ne0\), (G85.1) permits a basis over
\(\overline K\) in which
\[
 \rho(c)=\begin{pmatrix}1&0\\0&-1\end{pmatrix}.
\]
Put \(e_+=(1+\rho(c))/2\) and \(e_-=(1-\rho(c))/2\).
For every \(g\), the two diagonal entries of \(\rho(g)\) are
\[
 \frac{T(g)+T(cg)}2,\qquad
 \frac{T(g)-T(cg)}2,
 \tag{G85.2}
\]
and thus belong to \(K\).
Absolute irreducibility implies that some upper off-diagonal entry is
nonzero and that some lower off-diagonal entry is nonzero: otherwise
one of the two coordinate lines would be invariant.
Choose \(a\in\Gamma\) with nonzero upper entry and conjugate by a
diagonal matrix so that this entry becomes \(1\).
This conjugation leaves \(\rho(c)\) unchanged.
Writing the lower entry of \(\rho(g)\) as \(z_g\), we obtain
\[
 z_g=\operatorname{tr}\bigl(e_+\rho(a)e_-\rho(g)\bigr)\in K.
 \tag{G85.3}
\]
Indeed, expanding the two idempotents expresses this trace as a
\(K\)-linear combination of traces of group elements.
Choose \(a'\) with \(z_{a'}\ne0\).
If \(y_g\) is the upper entry of \(\rho(g)\), then
\[
 y_g z_{a'}=
 \operatorname{tr}\bigl(e_+\rho(g)e_-\rho(a')\bigr)\in K.
\]
Consequently \(y_g\in K\), proving the assertion.

For the last statement, at an absolutely irreducible specialization
choose four group elements whose matrices form a basis of
\(M_2\) over an algebraic closure of the specialization field.
The trace pairing on \(M_2\) is nondegenerate, so their trace Gram
determinant is nonzero there. The same determinant is an element of
\(S\), hence is nonzero generically. A reducible semisimple
two-dimensional generic representation is the sum of two characters,
and its matrix algebra has dimension at most two; its four-by-four
trace Gram determinants are therefore zero. This contradiction proves
generic absolute irreducibility.
Equivalently one may use Chenevier's
Definition--Proposition 2.18(v), whose trace-Gram criterion is exactly
this argument.
\end{proof}

\begin{proposition}[AC-G85-order: a finite free order containing the whole image]
\label{prop:G85-order}
In the setting of Lemma~\ref{lem:G85-split}, suppose that \(\Gamma\)
is profinite and \(T:\Gamma\to S\) is continuous.
There exists a unital \(\Lambda\)-algebra
\(\mathcal E\subset M_2(K)\), finite free of rank \(4D\) over
\(\Lambda\), with the following properties:
\[
 \rho(\Gamma)\subset\mathcal E^\times,\qquad
 S\mathcal E=\mathcal E,\qquad
 \operatorname{tr}(\mathcal E)\subset S.
 \tag{G85.4}
\]
Moreover \(\rho:\Gamma\to\mathcal E^\times\) is continuous.
No bound on a trace Gram denominator is required for the rank
or for the cardinality of a reduction of \(\mathcal E\).
\end{proposition}

\begin{proof}
\noindent\textbf{A finite order.}
The \(K\)-span of the matrices \(\rho(g)\) is \(M_2(K)\).
For completeness, it contains \(e_+\) and \(e_-\), and multiplication
by a matrix with nonzero upper, respectively lower, entry gives a
nonzero multiple of each off-diagonal matrix unit.
Together with \(e_+,e_-\) these span \(M_2(K)\).
Choose \(g_1,\ldots,g_4\) so that
\(v_i=\rho(g_i)\) is a \(K\)-basis, and put
\[
 G_{ij}=T(g_i g_j),\qquad d=\det(G)\in S\setminus\{0\}.
 \tag{G85.5}
\]
Nondegeneracy of the ordinary matrix trace pairing gives \(d\ne0\).
Solving for the coordinates of \(\rho(g)\) in this basis, using the
four numbers \(T(gg_j)\), shows
\[
 \rho(g)\in d^{-1}\sum_{i=1}^4 S v_i
 \quad\text{for every }g\in\Gamma.
\]
The \(S\)-module
\[
 \mathcal E_0=\sum_{g\in\Gamma}S\rho(g)
 \tag{G85.6}
\]
is consequently a submodule of a finite \(S\)-module.
It is finite because \(S\) is Noetherian.
It is a unital \(S\)-algebra since products of group matrices are
again group matrices, and \(\rho(g)^{-1}=\rho(g^{-1})\) belongs to it.
Its \(S\)-rank is four and its \(\Lambda\)-rank is \(4D\).
Its trace lies in \(S\).

\noindent\textbf{The reflexive order.}
Inside
\(V=\mathcal E_0\otimes_\Lambda K_0=M_2(K)\), define
\[
 \mathcal E=
 \bigcap_{\operatorname{ht}(\mathfrak p)=1}
       (\mathcal E_0)_{\mathfrak p},
 \tag{G85.7}
\]
where the primes belong to \(\operatorname{Spec}(\Lambda)\).
The height-one intersection description of reflexive hulls over a
Noetherian normal domain identifies this with
\(\mathcal E_0^{**}\); it is therefore finite and reflexive.
The ring \(\Lambda\) is regular local of dimension two, so finite
reflexive modules over it are free. Thus \(\mathcal E\) is free of
rank \(4D\).
These are precisely Stacks Project Tags 0AVC and 0B3N;
the latter follows from the depth-two criterion and Tag 00NT.

Each module \((\mathcal E_0)_{\mathfrak p}\) is also a unital algebra.
Their intersection is therefore a unital algebra, and it is stable
under multiplication by \(S\).
Furthermore
\[
 \operatorname{tr}(\mathcal E)
 \subset
 \bigcap_{\operatorname{ht}(\mathfrak p)=1} S_{\mathfrak p}
 =S.
 \tag{G85.8}
\]
The last equality holds because \(S\) is finite free, hence reflexive,
as a \(\Lambda\)-module. No assertion that the original order was
free is used.

\noindent\textbf{Continuity, including division by the fixed denominator.}
For \(g,h\in\Gamma\), the coordinate identity underlying (G85.5) gives
\[
 d\bigl(\rho(g)-\rho(h)\bigr)
 =
 \sum_i\bigl[\operatorname{adj}(G)
    (T(gg_j)-T(hg_j))_j\bigr]_i v_i.
 \tag{G85.9}
\]
Hence, if the four trace differences lie in
\(\mathfrak m^N S\), the left side lies in
\(\mathfrak m^N\mathcal E\).
Multiplication by \(d\) is an injective \(\Lambda\)-linear
endomorphism of the free module \(\mathcal E\); its determinant
\(\Delta\in\Lambda\) is nonzero.
Applying its adjugate to (G85.9) gives
\[
 \Delta\bigl(\rho(g)-\rho(h)\bigr)
       \in\mathfrak m^N\mathcal E.
 \tag{G85.10}
\]
The associated graded ring
\(\operatorname{gr}_{\mathfrak m}(\Lambda)\) is
\(k[X,Y]\), where \(k=\mathcal O/(\varpi)\), and is a domain.
Consequently
\(\operatorname{ord}_{\mathfrak m}(\Delta a)
=\operatorname{ord}_{\mathfrak m}(\Delta)
 +\operatorname{ord}_{\mathfrak m}(a)\)
for nonzero \(a\in\Lambda\).
This identity follows directly by multiplying the nonzero initial
homogeneous forms.
Apply it to the coordinates in a free basis of \(\mathcal E\).
For \(c_0=\operatorname{ord}_{\mathfrak m}(\Delta)\),
(G85.10) with \(N=n+c_0\) implies
\(\rho(g)-\rho(h)\in\mathfrak m^n\mathcal E\).
Continuity of the four functions \(T(gg_j)\) proves continuity of
\(\rho\).
The fixed denominator affects this proof of continuity only; it does
not change the rank \(4D\).
\end{proof}

\begin{theorem}[AC-G85: an explicit degree bound for the trace quotient]
\label{thm:G85}
Under the hypotheses above set
\[
 I=(p^4,W^4)\Lambda.
\]
There is an actual finite continuous group quotient \(Q_T\) of
\(\Gamma\) such that
\[
 |Q_T|\le p^{64Db},
 \tag{G85.11}
\]
and equality of two images in \(Q_T\) implies
\[
 T(g)-T(h)\in IS.
 \tag{G85.12}
\]
The same conclusion holds when the images are conjugate.

If \(\Gamma=G_{\mathbb Q,\Sigma}\), this quotient cuts out a finite
Galois extension \(L_T/\mathbb Q\), unramified outside \(\Sigma\),
of degree at most \(p^{64Db}\).
For any further fixed finite Galois extension \(L_0/\mathbb Q\),
\[
 [L_TL_0:\mathbb Q]\le
 [L_0:\mathbb Q]\,p^{64Db}.
 \tag{G85.13}
\]
Thus the trace congruence required by AC-G84 can be imposed together
with the fixed host, cubic character, cyclotomic modulo \(p^4\),
and AC-G83 unit-Kummer conditions in a finite Galois extension
with the displayed degree control.
\end{theorem}

\begin{proof}
Reduce the continuous representation of
Proposition~\ref{prop:G85-order} modulo \(I\mathcal E\), and take
\[
 Q_T=\operatorname{im}\bigl(
       \Gamma\longrightarrow(\mathcal E/I\mathcal E)^\times
                          \bigr).
\]
This is a group quotient, rather than merely the finite set of
values of a trace.
The ideal \(I\) is open in \(\Lambda\): if
\(p\mathcal O=\varpi^{e_0}\mathcal O\), then
\(\mathfrak m^{4e_0+3}\subset I\).
Thus its kernel is open by the continuity already proved.

The ring \(\Lambda/I\) has four coefficients in
\(\mathcal O/p^4\mathcal O\), so
\[
 |\Lambda/I|=p^{16b},\qquad
 |\mathcal E/I\mathcal E|
     =|\Lambda/I|^{4D}=p^{64Db}.
\]
This gives (G85.11).
The trace map carries \(I\mathcal E\) into \(IS\), proving
(G85.12). For conjugate images in the image group \(Q_T\), lift the
conjugator to \(\rho(\Gamma)\) and apply cyclicity of matrix trace.
If \(T\) was initially written on \(G_{\mathbb Q}\) but factors
through \(G_{\mathbb Q,\Sigma}\), the generic representation does too:
for any element \(u\) of the kernel, one has
\(\operatorname{tr}((\rho(u)-1)\rho(g))=T(ug)-T(g)=0\)
for every \(g\), and the group matrices span \(M_2(K)\).
Nondegeneracy of the trace pairing therefore gives \(\rho(u)=1\).
This proves the ramification assertion, not merely a statement about
traces of individual inertia elements.

The compositum degree inequality is elementary.
The actual element constructed in AC-G83 supplies a conjugacy class
in the image of the product of these quotients, so their simultaneous
prescription has no additional compatibility premise.
At a prime in this class, (G85.12) holds in \(S\), which is exactly
the coefficient ring on which the norm argument of AC-G84 is made.
\end{proof}

\paragraph{Application to the actual branch.}
The normalized degree-two determinant in AC-G78 is continuous and
factors through the Galois group unramified outside its fixed tame
level and \(p\).
Its determinant is the fixed odd character \(\eta_K\);
hence complex conjugation satisfies (G85.1).
The starting classical cuspidal modular representation is absolutely
irreducible. Lemma~\ref{lem:G85-split} therefore supplies generic
absolute irreducibility without any hypothesis of residual
irreducibility.
Normalize the finite branch over \(\Lambda\).
Excellence gives a finite normalization, and a normal
two-dimensional local ring is Cohen--Macaulay.
The normalization is finite free over the regular local ring
\(\Lambda\), as in AC-G84.
The inclusion of the original branch \(R\) into its normalization
\(S\) is continuous: both are finite \(\Lambda\)-modules and their
finite-module adic topologies are the ones used here.
Thus every assumption of Theorem~\ref{thm:G85} is met.
Here \(D\) is the rank of this actual normalized branch and
\(b=[E_0:\mathbb Q_p]\); neither is silently replaced by a constant
uniform in the original source parameters.
The congruence is in \(IS\).  The normalization version of AC-G84
requires only \(b_q-f\in IS\), produces an \(S\)-point, and restricts
that point to the original \(R\) before the actual Hecke-image step.
No equality \(IS\cap R=IR\) is asserted.

\paragraph{Source and proof scope.}
Chenevier's author text
\path{https://gaetan.chenevier.perso.math.cnrs.fr/articles/determinants.pdf},
Theorem 2.12 and Definition--Proposition 2.18
(printed pp.\ 28--33), supplies the field-valued determinant
realization and the trace-Gram criterion.
Their statements and printed proofs were inspected; their foundational
noncommutative-algebra inputs are retained, not re-proved here.
The split model, bounded order, reflexive-algebra construction,
continuity estimate, and quotient cardinality above are proved in full.
For the commutative-algebra inputs, the exact inspected Stacks Project
locators are Tags 0AVC, 0AVB, 0B3N, 00NT, 07QW, 035S and 0B3D.
The dimension and regular-sequence checks also use Tags 00H8 and 00N6,
whose complete short proofs were inspected.
The use of Pan's actual normalized determinant and the original
branch is the already checked AC-G78 interface; unavailable executor
access does not count as a new replay of the supplied Lean files.

\paragraph{Geometric dictionary and remaining gate.}
Using the already audited SGA~1, Expose~V, \S\S4 and~7 framework,
the finite continuous quotient on \(G_{\mathbb Q,\Sigma}\)
is a finite etale Galois cover of
\(\operatorname{Spec}\mathbb Z[1/\Sigma]\).
The reflexive order is an integral model over the two-dimensional
weight base; no geometric motive or additional purity assertion is
deduced from it.
The theorem controls the degree of the finite extension in which the
auxiliary Frobenius condition is imposed.
It does not assert a small least prime, a bound within the original
counting allowance, or the source's per-curve \(J\)-containment.
In particular the supplied theorem A is not substituted for
nonabelian GRH.
